\documentclass[../DoAn.tex]{subfiles}
\begin{document}

\begin{center}
    \Large{\textbf{TÓM TẮT NỘI DUNG ĐỒ ÁN}}\\
\end{center}
% \vspace{0.5cm}

% \noindent\textbf{Tên đề tài:} Trợ lý ảo luật lao động Việt Nam

\vspace{0.8cm}

Đồ án xây dựng hệ thống \textbf{hỏi đáp pháp luật lao động có trích dẫn} dựa trên Microsoft \gls{GraphRAG}~\cite{graphrag2024}, nhằm khắc phục hạn chế của tra cứu từ khóa truyền thống~\cite{robertson1994bm25} (không suy luận chéo văn bản) và chatbot \gls{LLM} tổng quát (dễ hallucination~\cite{ji2023hallucination}, trích dẫn sai lỗi thời).

\textbf{Đóng góp chính} là thiết kế \textbf{ontology hai lớp với canonicalization}: Lớp 1 (deterministic) biểu diễn cấu trúc pháp lý xác định không dùng LLM, đảm bảo 100\% nhất quán cho trích dẫn Điều/Khoản; Lớp 2 (LLM extraction) biểu diễn ngữ nghĩa domain với linh hoạt. Từ thiết kế này, hệ thống mở rộng sang suy luận đa bước (duyệt quan hệ pháp lý thay vì vector similarity), lọc trích dẫn hết hiệu lực, và phân tích tuân thủ hợp đồng kết hợp rule xác định với graph grounding.

Kết quả cho thấy: phương pháp đề xuất vượt trội LLM thuần và RAG vector truyền thống trên các câu hỏi yêu cầu trích dẫn chính xác và suy luận chéo quý phạm. Hệ thống đạt khả năng suy luận 100\% trên nhóm câu hỏi multi-hop (so với 33\% của Local Search đơn thuần), minh chứng rằng cấu trúc pháp lý xác định + ontology hai lớp là nền tảng cần thiết cho hỏi đáp pháp lý đáng tin cậy.

\textbf{Hạn chế}: corpus hẹp (lao động), bộ đánh giá nhỏ chưa qua chuyên gia. Tuy nhiên, pattern thiết kế có thể mở rộng sang các domain pháp lý khác và tái sử dụng trong các ứng dụng Legal AI tương tự. Hệ thống mang tính hỗ trợ, không thay thế tư vấn pháp lý chính thức.

\vspace{1.5cm}

\begin{flushright}
\textbf{Sinh viên thực hiện}\\[0.1cm]
\textit{(Ký và ghi rõ họ tên)}\\[2.5cm]
\textbf{Nguyễn Đạt Trọng}
\end{flushright}

\end{document}
